\documentclass[../cuaderno6.tex]{subfiles}

\begin{document}

\subsection{Polvo}
En el caso del polvo, resulta notorio lo cercano de los resultados obtenidos, los casos del campo coercitivo y la magnetización remanente resultan análogos al visible en la figura \ref{fig:M_max_segun_H_a_Polvo_Fa0_24_09} correspondiente a la magnetización máxima. Se considera que debido a su forma, polvo densamente empaquetado en un contenedor cilíndrico, rotaciones no deberían alterar el resultado de las mediciones.
\begin{figure}[H]
    \centering
    \includegraphics[width=0.5\textwidth]{../img/Fa0/Polvo/M_max_segun_H_a_Polvo_Fa0_24_09.png}
    \caption{Magnetización máxima $M_{max}$ en función del campo aplicado $H_{max}$.}
    \label{fig:M_max_segun_H_a_Polvo_Fa0_24_09}
\end{figure}

En cuanto al resultado del ajuste para el campo coercitivo $H_c^0$, obtenemos una mínima discrepancia, también vista en el factor de desmagnetización $N_d$, respectivamente en las figuras \ref{fig:H0_c_angulo_minors_Polvo} y \ref{fig:N_d_angulo_major_Polvo.png}

\begin{figure}[H]
    \centering
    \includegraphics[width=0.5\textwidth]{../img/Fa0/Polvo/H0_c_angulo_minors_Polvo.png}
    \caption{Gráfico de la pendiente del ajuste $H_c^0$ en función del ángulo.}
    \label{fig:H0_c_angulo_minors_Polvo}
\end{figure}
\begin{figure}[H]
    \centering
    \includegraphics[width=0.5\textwidth]{../../img/Fa0/Polvo/N_d_angulo_major_Polvo.png}

    \caption{Factor de desmagnetización $N_d$ en función del ángulo.}
    \label{fig:N_d_angulo_major_Polvo.png}
\end{figure}


\end{document}